\subsection{先选事件，再选参考系}
洛伦兹变换变换的是同一事件的坐标。解题时先为事件命名，再列出哪一系中哪个坐标差为零。常用的三个约束不能相互替代：匀速钟的两次读数在该钟的静止惯性系满足 $\Delta x'=0$；运动杆在实验室测长要求 $\Delta t=0$；杆的固有长度要求在杆系同时测两端。
若 $\Delta x=600\,\mathrm m$、$\Delta t=1.00\,\mu\mathrm s$，取 $c\simeq3.00\times10^8\,\mathrm{m/s}$，则 $c\Delta t=300\,\mathrm m<\Delta x$，两事件类空。
使两事件同时的换系速度由 $\Delta t'=0$ 得到
\[
 v=\frac{c^2\Delta t}{\Delta x}=0.50c.
\]
再改用 $v=0.80c$，有
\[
 \Delta t'=\frac53(1.00-1.60)\,\mu\mathrm s=-1.00\,\mu\mathrm s.
\]
事件先后可以反转，但两事件没有亚光速因果联系，因此不构成因果矛盾。反过来，若两事件类时，使它们同地的速度为 $v=\Delta x/\Delta t$，其绝对值小于 $c$；此时才存在连接两事件的钟的静止系。
